%% Un paso de tiempo del esquema implícito (Alya nastin): BDF2 + iteración no lineal +
%% Orthomin(1) que conserva la continuidad (Houzeaux et al. 2011, Algoritmo 5) + subescalas dinámicas.
\begin{tikzpicture}[
    caja/.style={fill=#1!12,draw=tinta,line width=0.6pt,rounded corners=3pt,align=left,inner sep=5pt,font=\scriptsize},
    sol/.style={fill=#1,text=white,rounded corners=1.5pt,inner sep=1.4pt,font=\tiny\bfseries,minimum height=0pt,minimum width=0pt,text width={},align=center},
    flecha/.style={-{Stealth[length=5pt]},tinta,line width=0.8pt}]
  \newcommand\mom{\tikz[baseline=(c.base)]\node[sol=ingEsq](c){MOMENTUM};}
  \newcommand\con{\tikz[baseline=(c.base)]\node[sol=ingTau](c){CONTINUITY};}
  % --- las tres piezas de una iteración no lineal ---
  \node[caja=ingEst,text width=2.3cm,minimum height=3.3cm] (ens) at (1.95,0)
    {\textbf{1 · Ensamblar}\\[3pt]
     Galerkin + VMS\\ ($\tau_1$, $\tau_2$)\\[3pt]
     BDF2 en $t$\\ {\color{tenue}(Euler los 20 primeros pasos)}\\[3pt]
     {\color{tenue}OSS: proyección de la it.\ anterior}};
  \node[caja=ingEsq,text width=5.2cm,minimum height=3.3cm,right=0.4cm of ens] (sch)
    {\textbf{2 · Schur: Orthomin(1)}, conserva la continuidad\\[4pt]
     \begin{tabular}{@{}l@{\hspace{4pt}}l@{}}
       \mom & $A_{uu}\,\boldsymbol u = \boldsymbol b_u - A_{up}\,p$\\[2pt]
       \con & $Q\,z = r$\\[2pt]
       \mom & $A_{uu}\,\boldsymbol v = A_{up}\,z$\\[2pt]
             & $\alpha$ óptimo; actualiza $\boldsymbol u$, $p$\\[2pt]
       \con & $Q\,z = r^{+}$; corrige $\boldsymbol u$
     \end{tabular}\\[3pt]
     {\color{tenue}\tiny GMRES + DIAGONAL · DEFLATED\_CG + LINELET · PRECONDITIONER: DT}};
  \node[caja=ingSub,text width=2.3cm,minimum height=3.3cm,right=0.4cm of sch] (sub)
    {\textbf{3 · Subescalas}\\[3pt]
     {\footnotesize$\partial_t\tilde{\boldsymbol u} + \tau_1^{-1}\tilde{\boldsymbol u}$}\\ {\footnotesize$= \mathcal P(\mathcal R_u)$}\\[3pt]
     Newton por elemento\\ {\color{tenue}tol $10^{-5}$}\\[3pt]
     {\color{tenue}solo en casos \caso{din}}};
  \draw[flecha] (ens) -- (sch); \draw[flecha] (sch) -- (sub);
  % --- marco de la iteración no lineal y su vuelta ---
  \begin{scope}[on background layer]
    \node[fill=tinta!4,rounded corners=5pt,inner xsep=8pt,inner ysep=10pt,fit=(ens)(sch)(sub)] (it) {};
  \end{scope}
  \node[anchor=south west,font=\scriptsize\bfseries,inner sep=1pt] at ($(it.north west)+(0.15,0.02)$) {Iteración no lineal};
  \draw[flecha,ingEst] (sub.south) -- ++(0,-0.62) -| node[pos=0.25,below,font=\scriptsize,text=tinta]
    {¿residuo $> 10^{-3}$? otra iteración (máx.\ 10)} (ens.south);
  % --- entrada y salida del paso ---
  \draw[flecha] ($(it.west)+(-0.65,0)$) node[left,font=\small]{$n$} -- (it.west);
  \draw[flecha] (it.east) -- ++(0.55,0) node[right,font=\small]{$n{+}1$};
\end{tikzpicture}
